\documentclass[10pt,a4paper]{amsart}
\usepackage[margin=19mm]{geometry}
\usepackage{amsmath,amssymb,mathtools}
\usepackage[scr=rsfs,cal=euler]{mathalfa}
\usepackage{xfrac}
\usepackage[unicode,hidelinks,bookmarks=false]{hyperref}
\newcommand{\ie}{i.e.\ }
\newcommand{\cf}{cf.\ }
\newcommand{\resp}{respectively\ }
\newcommand{\Subsection}{\subsection}
\providecommand{\nobreakdash}{\nobreak\hbox{-}}
\setlength{\emergencystretch}{2em}
\multlinegap0pt
\usepackage{enumerate}
\usepackage{mathtools}
\usepackage{amssymb}
\usepackage{bm}
\usepackage{tikz}
\usepackage{xcolor}
\usepackage{float}
\newtheorem{proposition}{Proposition}
\DeclareMathOperator{\CC}{CC}
\def\D{{\mathbb D}}
\def\R{\mathbb{R}}
\def\cM{{\mathcal M}}
\def\rd{{\rm d}}
\title{Contact $(+1)$-surgeries and algebraic overtwistedness}
\author{Zhengyi Zhou}
\date{Source: arXiv:2307.12635v3 (CC BY 4.0)}
\begin{document}
\maketitle

\noindent\emph{Source and licence.} Contact $(+1)$-surgeries and algebraic overtwistedness. Version arXiv:2307.12635v3 (CC BY 4.0), distributed by arXiv under the Creative Commons Attribution 4.0 International licence, http://creativecommons.org/licenses/by/4.0/, as recorded in the arXiv licence metadata for this exact version and shown on the abstract page https://arxiv.org/abs/2307.12635v3. Full licence notice: \texttt{licenses/arxiv-2307.12635-CC-BY-4.0.txt}.\par\smallskip

\noindent\textit{Editorial scope.} Counted unit: the article's proposition prop:area (source0.tex lines 849--855). The article's complete original proof of it is delivered with it (source0.tex lines 856--862, one \texttt{proof} environment of 7 lines). No mathematics is written, repaired or re-derived: the statement and the proof are the article's own lines, in the article's own notation, and no retained line is substituted or re-flowed.  Retained uncounted as the source-local context the proof needs: lines 799--799; lines 808--810; lines 837--839; lines 844--846.  Nothing was omitted from any retained span, so the package carries no omission record.  No retained line contains a citation, so the wrapper carries no bibliography and no bibliography entry.  No retained line contains a figure environment, an image-inclusion command or drawing code, so no image is delivered and the package declares no asset dependency.  Editorial formatting only, in addition: an AMS \texttt{amsart} preamble that restates the article's own declarations and macros used by the retained lines, namely the environments \texttt{proposition} on separate counters, together with the standard packages those lines need.  The retained environments print in this document as Proposition 1 in order of appearance, whereas the article prints the same items with the numbering of its own full document.  The complete native source of the article as distributed in the arXiv e-print for this exact version is bundled unchanged as \texttt{sources/144696/source0.tex}.
\subsubsection{The first case: $c_1(Y_1),c_1(Y_2)$ are torsion}\label{sss:421} We can assume that the SFT degrees of orbits with periods up to $8$ are strictly positive; these degrees are globally defined. So even though $\CC(Y_1\cup Y_2,\lambda|_{\partial_-W})$ may not have an augmentation, the part with action at most $8$ does. Specifically, we have the trivial map $\epsilon_0$ sending every generator to zero (by our degree assumptions), or equivalently the augmentation induced from the natural filling of $Y_1$ that is trivial on orbits of $Y_2$ (which are invisible with the action threshold above). We consider a Hamiltonian $H$ on the completion of $(Y_1\#Y_2, 2\lambda|_{\partial_+W})$ with slope $4$ and a Hamiltonian $G$ on the completion of $(Y_1\cup Y_2,\frac{1}{2}\lambda|_{\partial_-W})$ with slope $4$. Using $H_{VT}$ (and the cascades model to deal with the $S^1$ family of Hamiltonian orbits) we have the following Viterbo transfer map

\begin{equation}\label{eqn:vit}
    \phi_{Viterbo}(\overline{\gamma}_{\#})=\overline{\gamma}+\sum c_i\overline{\beta}_i
\end{equation}

\begin{equation}\label{eqn:area}
    \int u^*\rd\lambda \ge \int (\pi_{\D}u)^*\rd (\frac{r^2}{2\pi}\rd \theta)\ge 0,
\end{equation}

\begin{equation}\label{eqn:condition}
    \frac{1}{f(p)}-1<\int \gamma_{h,i}^*\lambda
\end{equation}

\begin{proposition}\label{prop:area}
    With the setup above, for $Y=Y_1\#Y_2$ and $\overline{W}$ the connected sum cobordism as in \S \ref{sss:421}, we have the following:
    \begin{enumerate}
        \item\label{a1} $\overline{\cM}_{Y,H}(\overline{\gamma}_p,\overline{\gamma}_{h,i}^k)=\emptyset$, $\overline{\cM}_{\overline{W},H}(\overline{\gamma}_p,\overline{\gamma}_{h,i}^k)=\emptyset$ for any $k$, where $p$ is not the minimum point of $f$.
        \item\label{a2}  The Viterbo transfer map in \eqref{eqn:vit} maps $\overline{\gamma}_{h,i}^k$ to zero.
    \end{enumerate}
\end{proposition}

\begin{proof}
     Since the linking number of $\gamma_p$ around the binding $\partial V\times\{0\}$ is $1$ and the linking number of $\gamma_h^k$ around $\partial V\times \{0\}$ is zero, for any curve $u$ in $\overline{\cM}_{Y,H}(\overline{\gamma}_p,\overline{\gamma}_{h,i}^k,\Gamma)$, we must have 
     $$\int \gamma_p^*\lambda-\int(\gamma_{h,i}^k)^*\lambda\ge \int_{u^{-1}(\partial V \times \D \times \R_t)}u^*\rd \lambda\ge \int_{u^{-1}(\partial V \times \D \times \R_t)}(\pi_{\D}u)^*(\frac{r^2}{2\pi}\rd \theta)=1,$$
     where the second inequality is \eqref{eqn:area} and the last equality follows from the linking number. Hence we arrive at a contradiction with \eqref{eqn:condition}. The same applies to $\overline{W}$ as the connected sum was applied outside of the $\partial V\times \D$ region.

    The intersection number of the curve $u$ in the Viterbo transfer map from $\overline{\gamma}_{h,i}^k$ with $\partial V \times \{0\}$ is negative, as orbits on $Y_1$ have positive linking with $\partial V$. However, this contradicts that $\pi_{\D}u$ is holomorphic, which implies that the intersection number is non-negative.  
\end{proof}

\end{document}
